\documentclass{scrartcl}
\usepackage[utf8]{inputenc}
\usepackage[T1]{fontenc}
\usepackage{geometry}
\usepackage{graphicx}
\usepackage{float}
\usepackage{hyperref}
\usepackage[english,indonesian]{babel}

\title{Optimasi Pemrosesan Template LaTeX Menggunakan Pendekatan Berbasis AI}
\author{Academic Support Agent}
\date{\today}

\begin{document}
\maketitle

\begin{abstract}

\end{abstract}

\chapter{PENDAHULUAN}


\chapter{TINJAUAN PUSTAKA}


% Original content (commented out):
% \documentclass[a4paper, 12pt, DIV=12]{scrartcl}
\usepackage[utf8]{inputenc}
\usepackage[T1]{fontenc}
\usepackage[french]{babel}
\usepackage{fancyhdr}
\usepackage{lastpage}

\usepackage{blindtext}
\usepackage{tikz}
\usepackage{graphicx}
\usepackage{tikzpagenodes}
\usepackage{tabularx}
\graphicspath{ {./images/} }
\usepackage{hyperref}
\usepackage{biblatex}


\title{Projet P5}
\author{auteur}
\date{September 2021}

\begin{document}


\pagestyle{fancy}
\fancyhf{}
\fancyhead[L]{Titre du projet} %en-tête de gauche
\fancyhead[R]{Rapport de projet PX/TB} %en-tête de droite
\fancyfoot[L]{auteur} %pied de page gauche
\fancyfoot[R]{\thepage/\pageref{LastPage}} %pied de page droite, numéro de page/total de page

\renewcommand{\headrulewidth}{1pt} %taille des lignes de séparation en-tête et pied de page
\renewcommand{\footrulewidth}{1pt}

\begin{titlepage}
\thispagestyle{empty}

\begin{tikzpicture}[remember picture,overlay] %logos, north west = haut gauche (comme sur une carte) même chose pour les autres coins
 \node[anchor=north west] at ([yshift=1.4cm]current page text area.north west) {\includegraphics[height=0.5in]{images/logo_heia-fr.jpg}};
 \node[anchor=south east] at ([yshift=-2.1cm]current page text area.south east) {\includegraphics[height=0.5in]{images/logo_heia-fr_version_courte.jpg}};
\end{tikzpicture}


   \begin{center}
    \vspace{0.5cm}
    \Large
   Année académique 21/22\\
   Rapport de projet PX\\
   \vspace{1cm}
   \LARGE
   Titre du projet\\
   \vspace{1cm}
   \includegraphics[height=8cm]{images/image_canevas.png}\\
   \vspace{1cm}
   \normalsize
   \begin{tabularx}{\textwidth} { %tableau pour créer 2 colonnes 
   >{\raggedright\arraybackslash}X 
   >{\raggedright\arraybackslash}X  }
        \textbf{Option} & Option\\
        & \\
        \textbf{Etudiant} & nom\_etudiant\\
        & \\
        \textbf{Enseignants responsables} & profs responsables\\
        & \\
        \textbf{Mendant} & Mendant \\ & rue, numéro\\ & code postal, ville \\
        & \\
        \textbf{Numéro de projet} & GM\_2122\_PX\_XX\_elève\_PROF1\_PROF2 \\
        & \\
        \textbf{Date} &  A METTRE A JOUR 
    \end{tabularx}
   \end{center}
\end{titlepage}


\renewcommand\thesection{\roman{section}}
\section{remerciements}
\blindtext
\newpage

\section{Résumé}
\blindtext
\newpage

\section{Abstract}
\blindtext
\newpage

\section{Table des matières}
\tableofcontents
\newpage

\section{Nomenclature}
\newpage

\setcounter{section}{0}
\renewcommand\thesection{\arabic{section}}


\section{Introduction}

\blindtext[2]
\newpage

\section{But du travail}
\newpage

\section{Documentation}
\newpage

\section{Description de l'exécution}
\newpage


\section{Résultats et analyse}
\newpage

\section{Analyse financière}
\newpage

\section{Conclusions}
\newpage

\section{Déclaration d'honneur}
\newpage

\section{Bibliographie}
\printbibliography{}
\newpage

\section*{Annexes}
\addcontentsline{toc}{section}{Annexes}


\end{document}

\end{document}